\section{Dual-Dimension Taxonomy}
\label{sec:taxonomy}

Section~\ref{sec:unified-theoretical-framework} unifies modern LLM optimizers
from two complementary views.  The universal meta-pipeline asks where an
optimizer intervenes in a single update, whereas the LMO-driven four-axis
decomposition asks how the intervention changes the update direction,
curvature estimate, gradient estimator, and state representation.  This
section turns those two views into a survey taxonomy that can also guide the
benchmark design.  The goal is not to list optimizers chronologically or to
assign informal empirical labels.  The goal is to construct a mechanism
coordinate system that supports literature organization, fair experimental
grouping, and interpretable trade-off analysis.

The taxonomy has two dimensions.  \emph{Dimension A} is a methodological
taxonomy: each optimizer receives one primary mechanism label so that
Section~\ref{sec:optimizer-families} can present T1--T5 as non-overlapping
families.  \emph{Dimension B} is an effect-objective taxonomy: each optimizer
may receive multiple objective labels describing the training properties it
claims to improve or should be evaluated against, such as convergence,
per-step cost, memory, stability, hyperparameter robustness, and
generalization.  The two dimensions answer different questions: what does the
optimizer change, and what outcome is the change intended to improve?

\subsection{Taxonomy Design Principles}
\label{sec:taxonomy-principles}

A single tree taxonomy is not expressive enough for modern LLM optimizers. Lion, for example, can reasonably be described as an Adam alternative, a sign-based optimizer, and a memory-friendly optimizer. All three descriptions are defensible: Lion keeps a first-order momentum structure, replaces continuous-valued directions with a sign map, and removes Adam second-moment state. However, the labels answer different questions.
The sign map is the primary update mechanism, whereas memory saving is an effect of that mechanism.  If Lion is placed only under ``memory-efficient optimizers,'' it is grouped with AdaFactor, 8-bit Adam, Adam-mini, and LOMO, although those methods mainly alter state storage, state sharing, or gradient lifetime. If Lion is copied into every relevant branch, the taxonomy loses its role as a clean survey organization and an experimental grouping rule.

The same ambiguity appears throughout the optimizer landscape.  GaLore
changes the subspace in which gradients and optimizer states live, whereas
Q-GaLore adds low-bit state quantization on top of that subspace mechanism.
SAM can wrap SGD or AdamW, but its defining contribution is a
sharpness-aware perturbation and second gradient evaluation.  LAMB first
performs an Adam-style element-wise update and then applies a layer-wise trust
ratio at writeback.  Cautious AdamW and Cautious Lion inherit different base
optimizers but share the same direction-consistency filter.  These examples
show why the taxonomy must separate primary mechanism, base optimizer,
composable wrapper, and intended effect.

We therefore use three design principles.

\textbf{First, Dimension A uses a single primary mechanism label.}
Each optimizer is assigned to exactly one main family. The assignment is not based on the name of the optimizer or chronology, but on the component whose removal would make the method collapse to a simpler baseline. Direct modifications of Adam-style moments, bias correction, time scales, iterate
averaging, or global step-size adaptation belong to T1.  Matrix routing,
spectral orthogonalization, Kronecker-factored preconditioning, and low-rank
subspace projection belong to T2.  Sign maps and related irreversible
direction discretization belong to T3.  State sharing, compression,
quantization, or elimination belong to T4.  Perturbations, explicit curvature
wrappers, post-update filters, and layer-wise trust-region rules belong to T5.
Composite methods are classified by their incremental contribution and record
secondary tags in the cross-dimension matrix.

\textbf{Second, family boundaries are aligned with the meta-pipeline.}
The five families are not arbitrary names.  T1 mainly modifies element-wise
state evolution in S3 and local scaling in S5.  T2 changes the update space
through S1--S4 matrix routing, matrix transformation, structured state, and
reconstruction.  T3 introduces irreversible direction discretization in S2
and often simplifies S3.  T4 acts on S3/S4 state representation and gradient
lifetime.  T5 acts near S5 through update writeback, post-processing, or
geometric regularization.  Given an update rule, locating its dominant
non-identity pipeline operation is usually enough to assign the method label.

\textbf{Third, Dimension B uses multi-label effect annotations.}
Effect objectives should not be mutually exclusive categories.  A method may
target convergence efficiency and hyperparameter robustness at the same time,
or it may reduce memory while hurting convergence or stability.  An effect
label means that the method claims to improve, is designed to improve, or must
be evaluated on that property.
This separation avoids equating ``designed to save memory'' with ``better
under a fixed hardware budget,'' and avoids assuming that every
sharpness-aware method improves generalization under every protocol.

\subsection{Methodological Taxonomy}
\label{sec:taxonomy-method}

Dimension A organizes the surveyed set of 108 optimizers into five families
and fifteen subclasses.  The five
families correspond to the dominant mechanism that transforms a stochastic
gradient into a parameter update.  The main text keeps the organizing
principles and family definitions that are needed for the survey chapters that
follow, with a compact family-level summary given in
Table~\ref{tab:taxonomy_cross_matrix}.

\textbf{T1: Element-wise adaptive moment and scalar control.}
T1 retains the Adam-style element-wise processing paradigm, but it is broader
than a strict set of Adam variants.  Its core members maintain first moments,
second moments, or related variance surrogates and normalize each parameter
independently.  Its boundary members are scalar first-order predecessors or
outer-control schemes that do not introduce matrix routing, sign
quantization, state compression, or sharpness-aware wrapper gradients.
Differences within the family come from temporal structure, variance
estimation, iterate averaging, corrected decay, and global step-size rules.
In the four-axis view, most T1 methods use the identity analysis basis, a
diagonal Fisher or variance estimate, square-root preconditioning for the
Adam-like core, and full state storage.

\textbf{T2: Matrix-level structural methods.}
T2 treats two-dimensional Transformer weights as matrices rather than as
flattened vectors.  The family includes spectral orthogonalization,
Kronecker-factored preconditioning, and low-rank subspace projection.  These
methods share nontrivial S1 routing and S2 transformation: they either
orthogonalize matrix momentum, maintain row/column curvature factors, or
project gradients into a low-dimensional subspace before updating state.  T2
is the family that most visibly changes Axis~I.  Muon-style methods use a
dynamic SVD or polar basis, Shampoo and SOAP use Kronecker or Fisher eigenbases,
and GaLore changes the effective optimizer-state space through projection.

\textbf{T3: Discretization and directional quantization.}
T3 is defined by sign maps or related direction-discretization operations.
The core idea is to keep coarse direction information while discarding or
quantizing magnitude.  Lion is the representative method: it forms a signed
direction from a gradient-momentum interpolation and keeps only first-order
momentum.  Compared with T1, T3 is not merely a cheaper second-moment
estimator.  It changes the geometry of the update direction.  In the LMO view,
the family is naturally associated with $\ell_\infty$-type direction
selection and weak or absent Hessian preconditioning.  Methods whose primary
increment is a binary or stochastic mask applied to an already computed
update are treated as T5.3 post-update filtering rather than as T3, even when
the mask is motivated by directional agreement.

\textbf{T4: State compression and structural aggregation.}
T4 directly targets the optimizer-state memory bottleneck.  It includes
row/column factorization, low-bit state quantization, block- or layer-level
state sharing, and fused backprop-update schemes.  Its distinction from T2.3
is the object being compressed.  T2.3 first changes the mathematical subspace
of the gradient signal and then runs an optimizer in that subspace.  T4 mainly
changes how the state itself is stored, shared, quantized, or released.  In
the four-axis view, T4 is the most direct instantiation of state compression in Axis~II.

\textbf{T5: Curvature-aware and geometric regularization.}
T5 does not primarily introduce a new element-wise moment estimator or matrix
transform.  Instead, it adds geometric constraints near update finalization.
Sharpness-aware methods change the point at which gradients are evaluated.
Hessian-guided methods replace or augment curvature signals and often add
trust-region clipping.  Post-processing methods centralize, project,
normalize, mask, or sparsify an already computed update.  Layer-wise
trust-region methods rescale updates according to parameter and update norms.
Because many T5 methods are wrappers around a base optimizer, the base
optimizer and the T5 mechanism must be recorded separately.

Boundary cases are resolved by the primary incremental mechanism.  Direct
Adam variants remain T1 even if they claim stability or generalization
benefits, because their main mechanism is still element-wise moment
estimation or step-size control.  Q-GaLore is T4.2 when the focus is the
additional quantized-state mechanism, while base GaLore remains T2.3.  Magma and MGUP are T5.3 because their primary operation is a selective
update mask.  SAM, Cautious optimizers, and
LAMB are classified by wrapper-level contributions in T5 rather than by their
possible AdamW or Lion base optimizers.  This keeps
the main survey families mutually exclusive while leaving room for secondary
tags in the appendix-level optimizer matrix.

\subsection{Objective-Oriented Taxonomy}
\label{sec:taxonomy-objective}

Dimension B describes what an optimizer is intended to improve.  We organize
effect objectives by evaluation cost rather than by method mechanism, because
the feasibility of a benchmark depends on how each objective is measured.  We
use six objectives: O1 convergence efficiency, O2 per-step computational
cost, O3 memory overhead, O4 training stability, O5 hyperparameter
robustness, and O6 generalization quality.  These labels are not mutually
exclusive and should not be read as intrinsic family rankings.  They are
evaluation axes for organizing experiments and interpreting trade-offs.

The six objectives form three measurement layers.

\textbf{Layer 1: directly logged metrics.}
O1--O3 can be obtained from a single training run, timers, profilers, or
analytic byte/FLOP models.  O1 measures how fast loss decreases or a target
quality is reached under a fixed step, token, wall-clock, or compute budget.
O2 measures additional per-step computation relative to AdamW or SGD,
including matrix multiplications, orthogonalization iterations, extra
forward/backward passes, synchronization, and quantization/dequantization.
O3 measures memory from parameters, gradients, optimizer states, temporary
matrices, projection bases, and quantization buffers.

\textbf{Layer 2: single-run derived metrics.}
O4 and the lightweight version of O6 require post-processing of standard
training logs.  O4 measures whether the trajectory is robust to loss spikes,
gradient-norm bursts, overflow, divergence, and incomplete runs.  The
lightweight O6 signal comes from validation loss, train-validation gap, or a
fixed validation protocol during pretraining.

\textbf{Layer 3: cross-configuration metrics.}
O5 and the full version of O6 require multiple training runs or downstream
evaluations.  O5 measures sensitivity to learning rate, weight decay,
momentum, warmup, batch size, and method-specific hyperparameters.  Full O6
measures downstream transfer, out-of-distribution retention, scale transfer,
or other quality metrics beyond pretraining loss.

\begin{table*}[!t]
  \centering
  \caption{Dimension-B effect objectives and measurement sources.}
  \label{tab:effect_objectives}
  \scriptsize
  \setlength{\tabcolsep}{1.5pt}
  \renewcommand{\arraystretch}{1.08}
  \begin{tabularx}{\textwidth}{@{}Y p{0.17\textwidth}p{0.14\textwidth}Y@{}}
    \toprule
    Definition & Data source & Extra cost & Typical outputs \\
    \midrule
    \rowcolor{famO1!18}\multicolumn{4}{c}{\textbf{\textcolor{famO1!62!black}{O1: Convergence Efficiency}}} \\
      Loss reduction and time-to-target under a fixed training budget
      & Train/validation loss logs
      & None
      & Final loss; steps-to-threshold; token efficiency \\
    \rowcolor{famO2!18}\multicolumn{4}{c}{\textbf{\textcolor{famO2!62!black}{O2: Step Cost}}} \\
      Extra per-step computation and synchronization relative to a baseline
      & Timers, FLOP analysis
      & Recorded during training
      & Step time; FLOPs; extra backward count \\
    \rowcolor{famO3!18}\multicolumn{4}{c}{\textbf{\textcolor{famO3!62!black}{O3: Memory}}} \\
      Memory from optimizer states and associated buffers
      & Memory profiler, byte model
      & Recorded during training
      & Peak memory; state and buffer bytes \\
    \rowcolor{famO4!18}\multicolumn{4}{c}{\textbf{\textcolor{famO4!62!black}{O4: Stability}}} \\
      Robustness to spikes, divergence, and gradient fluctuations
      & Loss and gradient-norm curves
      & Offline post-processing
      & Spike rate; gradient CV; divergence rate \\
    \rowcolor{famO5!18}\multicolumn{4}{c}{\textbf{\textcolor{famO5!62!black}{O5: Hyperparameter Robustness}}} \\
      Sensitivity to learning rate, decay, batch size, and other knobs
      & Multiple training runs
      & Search or transfer experiments
      & Usable LR interval; performance variance; tuning burden \\
    \rowcolor{famO6!18}\multicolumn{4}{c}{\textbf{\textcolor{famO6!62!black}{O6: Generalization}}} \\
      Quality beyond the training objective (validation, downstream, OOD, transfer)
      & Validation and downstream evaluation
      & Low for validation, high for full evaluation
      & Validation loss; generalization gap; downstream score \\
    \bottomrule
  \end{tabularx}
\end{table*}



Dimension B must remain neutral before benchmarking.  T2 methods may improve
O1 token efficiency through matrix structure, but their orthogonalization or
Kronecker operations may hurt O2.  T4 methods are designed for O3, but
aggressive compression may degrade O1 or O4.  T5 methods often target O4 and
O6, but SAM-style extra gradient evaluations impose a clear O2 cost.  Thus
effect labels are measurement prompts rather than substitutes for empirical
results.

For the family analyses in Section~\ref{sec:optimizer-families}, each
optimizer or subclass can be annotated by
\begin{equation}
  \operatorname{Effect}(A) \subseteq \{O1,O2,O3,O4,O5,O6\},
  \label{eq:effect-tags}
\end{equation}
where $A$ denotes an optimizer or subclass.  The annotation should distinguish
design intent from experimentally verified effect.  AdaFactor, for instance,
can be marked as O3-targeting under an analytic memory model, but whether it
also improves O1, O4, or O6 under a fixed training protocol must be decided
by the benchmark.

\subsection{Cross-Dimension Analysis}
\label{sec:taxonomy-cross}

The value of the dual taxonomy appears in cross-dimension analysis.
Dimension A gives mechanism labels, while Dimension B gives effect objectives.
Table~\ref{tab:taxonomy_cross_matrix} provides a compact matrix of the
objectives that each method family should emphasize in the benchmark.  The
matrix is a mechanism-informed prior, not a final empirical conclusion.

\begin{table}[!t]
  \centering
  \caption{Compact cross-matrix between method families and effect
  objectives.}
  \label{tab:taxonomy_cross_matrix}
  \small
  \newcommand{\ppos}{\textcolor{green!35!black}{\ensuremath{\boldsymbol{\Uparrow}}}}
  \newcommand{\pos}{\textcolor{green!35!black}{\ensuremath{\boldsymbol{\Uparrow}}}}
  \newcommand{\neu}{\tikz[baseline=-0.55ex]\draw[gray,line width=0.45pt] (0,0) circle (0.62ex);}
  \newcommand{\negv}{\textcolor{red!70!black}{\ensuremath{\boldsymbol{\Downarrow}}}}
  \newcommand{\posneu}{\textcolor{green!65!black}{\ensuremath{\uparrow}}}
  \newcommand{\neupos}{\textcolor{green!65!black}{\ensuremath{\uparrow}}}
  \newcommand{\neuneg}{\textcolor{red!55}{\ensuremath{\downarrow}}}
  \newcommand{\negneu}{\textcolor{red!55}{\ensuremath{\downarrow}}}
  \setlength{\tabcolsep}{5pt}
  \begin{tabularx}{\linewidth}{@{}>{\raggedright\arraybackslash}X*{6}{>{\centering\arraybackslash}m{0.065\linewidth}}@{}}
    \toprule
    Family & O1 & O2 & O3 & O4 & O5 & O6 \\
    \midrule
    T1 Element-wise adaptive moment and scalar control & \ppos & \neu & \negv & \pos & \pos & \neupos \\
    T2 Matrix-level structural methods & \ppos & \negv & \negneu & \pos & \neupos & \neupos \\
    T3 Discretized directions & \posneu & \pos & \pos & \posneu & \neu & \neu \\
    T4 State compression & \neuneg & \neupos & \ppos & \neuneg & \neu & \neuneg \\
    T5 Geometry regularization & \posneu & \negv & \neuneg & \ppos & \pos & \ppos \\
    \bottomrule
  \end{tabularx}
  \vspace{2pt}
  \begin{minipage}{0.98\linewidth}
    \footnotesize \emph{Note.} Dark double-line arrows denote strong priors:
    \ppos{} indicates a favorable target and \negv{} a likely cost. Pale
    single-line arrows denote conditional or protocol-sensitive favorable/cost
    effects. The hollow gray circle \neu{} denotes protocol-dependent
    neutrality. Arrows indicate objective favorability rather than raw metric
    direction.
  \end{minipage}
\end{table}


The symbols in Table~\ref{tab:taxonomy_cross_matrix} should not be read as a
global performance ranking.  A T2 method may be slower per step but still win
in wall-clock time if token efficiency improves enough.  A T4 method may look
worse in a fixed-small-model loss comparison while enabling a larger batch or
model under the same hardware budget.  A T5 method may improve generalization
only at particular data scales, training durations, or evaluation suites.

The cross view also identifies mechanisms that are likely to compose
orthogonally.  Low-rank projection and state quantization can be combined
because T2.3 changes the signal subspace, whereas T4.2 changes the state
representation.  Variance-reduced gradient estimates can be layered onto
AdamW, Lion, or Shampoo because variance reduction (Axis~II) is mostly independent of basis choice
and curvature representation.  Post-update filters can wrap AdamW or Lion
because T5.3 acts after the base direction has been computed.  Layer-wise
trust ratios can similarly wrap an element-wise adaptive update because they
primarily act in S5.

Not every cross-family combination is natural.  Mechanisms that occupy the
same meta-pipeline slot require an explicit ordering and a descent-quality
argument.  Spectral orthogonalization and low-rank projection are both strong
S2 constraints on a matrix signal, so combining them requires specifying whether
projection precedes orthogonalization or vice versa.  LOMO-style streaming
updates can conflict with methods that need global gradient statistics,
delayed basis updates, or a second gradient evaluation.  SAM-style methods
already require extra forward/backward computation, so adding expensive
Kronecker-factored preconditioning may be mathematically valid but practically
unattractive.  These constraints are not implementation details.  They are
composition boundaries that the taxonomy should make explicit.

The full appendix-level matrix should therefore record, for each optimizer,
three fields: the primary Dimension-A label, secondary mechanism tags, and
O1--O6 effect annotations.  The main text keeps only compact summaries such
as Table~\ref{tab:taxonomy_cross_matrix}; the family analyses instantiate the
matrix progressively by describing each subclass's meta-pipeline site,
four-axis coordinates, and expected benchmark signature.  In this way, the
taxonomy is a mechanism-aware
benchmark plan: a new optimizer entering the framework must specify which
pipeline stage it modifies, which four-axis coordinates it changes, and which
effect objectives it claims to improve.
